\documentclass[9pt,twocolumn]{extarticle}
\usepackage[letterpaper,textwidth=178mm,textheight=229mm,top=25mm,left=19mm,columnsep=6mm,noheadfoot]{geometry}
\usepackage{amsmath,amssymb,booktabs,array,url,times}
\usepackage[hidelinks]{hyperref}
\usepackage{titlesec}
\pagestyle{empty}
\setlength{\parindent}{1em}
\setlength{\parskip}{0pt}
\titleformat{\section}{\bfseries\centering\normalsize}{\thesection.}{0.6em}{\MakeUppercase}
\titlespacing*{\section}{0pt}{1.0ex plus .2ex minus .2ex}{0.6ex}
\titleformat{\subsection}{\bfseries\normalsize}{\thesubsection}{0.5em}{}
\titlespacing*{\subsection}{0pt}{0.8ex}{0.4ex}
\renewenvironment{abstract}{\begin{center}\bfseries ABSTRACT\end{center}\noindent}{\par}
\newenvironment{keywords}{\vspace{0.5em}\noindent\textbf{\textit{Index Terms}}---}{\par}
\title{\vspace*{2mm}\textbf{CAFA-IVR: COUNTERFACTUAL ASR FAILURE ATTRIBUTION\\FOR CONVERSATIONAL IVR TESTING}}
\author{Vijay Kumar Sridharan\\Independent Researcher\\Frisco, TX, USA\\\texttt{sridharanvijaykumar@gmail.com}}
\date{}
\begin{document}

\maketitle
\thispagestyle{empty}

\begin{abstract}
Automatic speech recognition (ASR) in conversational IVR is commonly evaluated with word error rate (WER), although user success depends on downstream intent and task behavior. We introduce CAFA-IVR, a vendor-neutral counterfactual testing framework that executes each semantic case through two paths: reference text to the downstream system, and audio to ASR to the same downstream system. A failure that appears only on the speech path is measured by ASR-Attributable Intent Failure Rate (ASR-IFR). CAFA-IVR also defines critical-entity scoring and a portable CI/CD contract. We validate the method with 540 locally executed ASR trials: 180 PocketSphinx decodes and 360 Conformer-CTC decodes across clean, telephony, noise, and speaking-rate conditions. In the neural study, clean ASR-IFR was 10.0\% and rose to 71.7\% at 5-dB noise. We additionally map the protocol to public human banking speech and release a reference scorer for reproducible adoption.
\end{abstract}

\begin{keywords}
automatic speech recognition, spoken language understanding, conversational IVR, failure attribution, reliability testing
\end{keywords}

\section{Introduction}
A conversational IVR is a layered system: audio capture and channel processing feed ASR, whose transcript is interpreted by natural-language understanding (NLU), an agent, or dialogue logic. When the observed intent or task is wrong, the final failure alone does not identify whether the speech layer or the downstream system caused it.

WER remains the standard lexical diagnostic,
\begin{equation}
\mathrm{WER}=\frac{S+D+I}{N},
\end{equation}
where $S$, $D$, and $I$ are substitutions, deletions, and insertions relative to $N$ reference words. However, lexical distance is not an attribution mechanism. Semantic-ASR research has similarly motivated alternatives or complements to pure lexical scoring \cite{kim2021,roy2021}, while spoken-language-understanding studies document downstream degradation under real-world input noise \cite{sengupta2021}. Bot test workbenches demonstrate the practicality of baseline-based intent and slot regression testing \cite{awslex}. The remaining operational question is: \emph{did the addition of the speech path cause an otherwise valid semantic test to fail?}

CAFA-IVR makes four contributions: (1) a paired text/audio counterfactual; (2) ASR-IFR; (3) a portable attribution and critical-entity contract suitable for CI/CD; and (4) reproducible real-decoding pilots plus an external human-banking-corpus mapping. The reference scorer, specification, examples, and empirical artifacts are openly available at \url{https://github.com/sridharanvijaykumar/cafa-ivr}.

\section{CAFA-IVR Framework}
For semantic case $i$, let $r_i$ be the reference text and $a_i$ its audio realization. CAFA-IVR executes
\begin{align}
T_i &: r_i \rightarrow \mathrm{Downstream},\\
A_i &: a_i \rightarrow \mathrm{ASR} \rightarrow \mathrm{Downstream}.
\end{align}
The downstream model, prompt, tools, dialogue state, and test oracle should be frozen between the two paths. Let $T_i=1$ and $A_i=1$ denote successful semantic outcomes. Table~\ref{tab:attrib} defines the attribution states.

\begin{figure}[t]
\centering
\fbox{\begin{minipage}{0.94\columnwidth}
\centering\footnotesize
Reference text $r_i$ $\longrightarrow$ Same downstream system $\longrightarrow T_i$\\[1mm]
Audio $a_i$ $\longrightarrow$ ASR $\longrightarrow$ Same downstream system $\longrightarrow A_i$\\[1mm]
$(T_i,A_i)$ $\longrightarrow$ Attribution state $\longrightarrow$ WER / CEER / ASR-IFR
\end{minipage}}
\caption{CAFA-IVR paired counterfactual execution. Only the audio path introduces ASR; downstream components and the task oracle remain frozen.}
\label{fig:cafa}
\end{figure}

\begin{table}[t]
\caption{CAFA-IVR attribution matrix.}
\label{tab:attrib}
\centering\scriptsize
\begin{tabular}{ccl}
\toprule
Text & Audio & Attribution\\
\midrule
Pass & Pass & HEALTHY\\
Pass & Fail & SPEECH\_ATTRIBUTABLE\\
Fail & Fail & DOWNSTREAM\_OR\_TEST\\
Fail & Pass & CONTEXT\_OR\_ORACLE\\
\bottomrule
\end{tabular}
\end{table}

Define $F_i=\mathbb{1}(T_i=1\land A_i=0)$. The primary metric is
\begin{equation}
\mathrm{ASR\mbox{-}IFR}=\frac{1}{N}\sum_{i=1}^{N}F_i.
\end{equation}
Thus ASR-IFR measures the fraction of the complete paired suite converted from a downstream-capable reference-text success to an audio-path failure. For teams that want risk conditioned only on eligible controls, CAFA-IVR additionally reports
\begin{equation}
\mathrm{cASR\mbox{-}IFR}=\frac{\sum_i F_i}{\sum_i T_i},
\end{equation}
while keeping the all-trial ASR-IFR as the primary release metric.

CAFA-IVR also retains WER and defines Critical-Entity Error Rate (CEER) over governed values such as amounts, dates, negation, authentication values, account types, and high-risk action verbs. If $C_i$ and $\hat C_i$ are normalized expected and observed critical-entity multisets,
\begin{equation}
\mathrm{CEER}=\frac{\sum_i (|C_i\setminus\hat C_i|+|\hat C_i\setminus C_i|)}{\sum_i |C_i|}.
\end{equation}
An optional consequence score (CIER) may be added, but severity weights are deployment policy rather than universal constants.

\section{Portable Test and Release Contract}
The minimum trial record contains a stable test ID, acoustic condition, reference text, ASR transcript, expected/observed semantic outcome, text-control result, audio-path task result, and engine/version metadata. Critical entities and risk tier are recommended for high-consequence journeys.

The framework is deliberately vendor-neutral. The downstream component can be a classical intent model, deterministic workflow, LLM agent, or tool-calling system, provided a reproducible task oracle exists. The operational algorithm is:

\begin{enumerate}
\item freeze candidate ASR and downstream versions;
\item run all reference-text controls;
\item run the identical cases through audio and ASR;
\item compute WER, ASR-IFR, and CEER where applicable;
\item compare with the last approved baseline;
\item route failures using Table~\ref{tab:attrib};
\item retain failed audio, transcript, versions, and outcome for replay.
\end{enumerate}

CAFA-IVR does not prescribe a universal numeric release threshold. Instead, teams can gate on regression relative to an approved baseline, with stricter policy for high-risk intents. For example, an organization may require no unexplained ASR-IFR increase and no new critical-entity errors on governed payment, fraud, or card-control journeys. WER remains a diagnostic but is not the sole release decision.

\section{Experimental Design}
\subsection{Legacy Pilot}
A 600-utterance benchmark spans 30 IVR intents. One utterance per intent was held out and the other 570 trained a TF-IDF/logistic-regression downstream classifier. Held-out text-control accuracy was 96.7\%. eSpeak generated audio under six deterministic conditions: clean, 8-kHz telephony, 15-dB noise, 5-dB noise, 1.15$\times$ rate, and 8-kHz plus 15-dB noise. PocketSphinx produced 180 actual ASR decodes. Overall WER was 88.9\%, intent accuracy 22.2\%, and ASR-IFR 75.0\%. This weak recognizer/synthetic-speech combination is retained as pipeline validation rather than a modern-ASR performance claim.

\subsection{Neural Pilot}
A second experiment used a 60-phrase closed IVR grammar (two phrases per intent) synthesized by ten eSpeak English voices. V01--V08 (480 clips) trained the acoustic model, V09 (60) was validation, and unseen V10 (60) was the test speaker. The acoustic model used 64-bin log-Mel features, a 96-dimensional representation, three Conformer blocks, four attention heads, a 256-dimensional feed-forward layer, kernel size 15, and character CTC output (516,029 parameters). Conformer combines convolutional local modeling with Transformer-style global modeling \cite{gulati2020}. The frozen epoch-20 model used a constrained 60-phrase CTC decoder. The downstream NLU was separately trained on the other 540 benchmark texts and obtained 80.0\% text-control accuracy on these 60 cases.

The V10 test audio was evaluated under the same six conditions, producing 360 actual neural decodes with no condition-specific retraining.

\section{Results}
Table~\ref{tab:studies} summarizes the two locally executed studies. They validate the same attribution pipeline under different recognizer assumptions and are not presented as an apples-to-apples ASR leaderboard.

\begin{table}[t]
\caption{Locally executed CAFA-IVR studies.}
\label{tab:studies}
\centering\scriptsize
\begin{tabular}{lrrrr}
\toprule
Study & Trials & Text ctrl. & WER & ASR-IFR\\
 & & Acc. (\%) & (\%) & (\%)\\
\midrule
PocketSphinx & 180 & 96.7 & 88.9 & 75.0\\
Conformer-CTC & 360 & 80.0 & 64.8 & 39.4\\
\bottomrule
\end{tabular}
\end{table}

Table~\ref{tab:neural} reports measured neural results. Clean WER was 11.5\%, exact phrase recognition 88.3\%, downstream intent accuracy 70.0\%, and ASR-IFR 10.0\%. Telephony increased ASR-IFR to 20.0\%, and severe 5-dB noise increased it to 71.7\%. WER can exceed 100\% when total word edits exceed the reference length.

\begin{table}[t]
\caption{Measured Conformer-CTC results, 60 trials/condition.}
\label{tab:neural}
\centering\scriptsize
\begin{tabular}{lrrrr}
\toprule
Condition & WER & Phrase & Intent & IFR\\
 & (\%) & Exact (\%) & Acc. (\%) & (\%)\\
\midrule
Clean & 11.5 & 88.3 & 70.0 & 10.0\\
8 kHz & 27.2 & 68.3 & 60.0 & 20.0\\
15 dB & 74.1 & 41.7 & 45.0 & 36.7\\
5 dB & 122.9 & 6.7 & 10.0 & 71.7\\
Rate 1.15$\times$ & 41.2 & 55.0 & 48.3 & 31.7\\
8 kHz + 15 dB & 112.1 & 18.3 & 13.3 & 66.7\\
\midrule
Overall & 64.8 & 46.4 & 41.1 & 39.4\\
\bottomrule
\end{tabular}
\end{table}

Nonparametric trial-level bootstrap intervals used 5,000 resamples with a fixed seed. Bootstrap 95\% intervals for clean speech were 4.2--20.1\% WER, 58.3--81.7\% intent accuracy, and 3.3--18.3\% ASR-IFR. Across all 360 trials, the corresponding intervals were 57.7--72.1\%, 35.8--46.1\%, and 34.4--44.4\%. The neural text control passed in 288 of 360 paired trials; 142 were speech-attributable failures, yielding 39.4\% all-trial ASR-IFR and 49.3\% conditional cASR-IFR.

The distinction from WER is visible at trial level. ``my physical card is missing'' became ``the card is missing'' (WER 0.40) while preserving \texttt{lost\_card}. By contrast, ``what are my latest payments'' became ``what are my last transactions'' (WER 0.40), changing \texttt{payment\_history} to \texttt{recent\_transactions}. Equal WER therefore need not imply equal semantic consequence.

The text counterfactual additionally prevented false ASR attribution. Twelve of the 60 neural reference texts failed the downstream NLU before speech was introduced. Across six conditions, 34 trials had an exact ASR transcript but belonged to those text-control failures; a voice-only dashboard could otherwise direct debugging toward ASR.


\section{Relation to Existing Evaluation Approaches}

CAFA-IVR differs from transcript-centric and downstream-only evaluation in the question it answers. Conventional WER measures lexical deviation between reference and ASR output. Semantic distance and Semantic-WER attempt to weight transcription errors by their semantic effect, while spoken-language-understanding robustness studies measure degradation of downstream intent or slot performance under noisy inputs \cite{kim2021,roy2021,sengupta2021}. These approaches are useful diagnostics, but they do not establish whether a particular downstream failure was introduced by the speech path or was already present in the downstream system.

CAFA-IVR instead uses a paired counterfactual control. For each semantic case, the reference text and the ASR transcript are evaluated against the same frozen downstream configuration and task oracle. Only a case that succeeds on the reference-text path and fails after speech is introduced is classified as speech-attributable. Thus, CAFA-IVR is not a replacement for WER or semantic transcript metrics; it adds failure attribution at the system boundary.

\begin{table}[t]
\caption{Evaluation focus of representative approaches.}
\label{tab:comparison}
\centering
\scriptsize
\begin{tabular}{lcccc}
\toprule
Approach & Lexical & Task & Text & Speech \\
 & error & outcome & control & attribution \\
\midrule
WER & \checkmark & -- & -- & -- \\
Semantic transcript metrics & \checkmark & partial & -- & -- \\
SLU robustness testing & partial & \checkmark & -- & -- \\
Voice-only end-to-end testing & optional & \checkmark & -- & -- \\
CAFA-IVR & \checkmark & \checkmark & \checkmark & \checkmark \\
\bottomrule
\end{tabular}
\end{table}

\subsection{Attribution Utility}
The paired control changes diagnosis even when transcript quality appears unproblematic. In the neural experiment, 12 of 60 reference texts already failed the downstream NLU before speech was introduced. Across the six acoustic conditions, 34 trials produced exact ASR transcripts while belonging to these text-control failures. A voice-only evaluation could therefore associate these cases with the speech pipeline despite the downstream failure being reproducible without ASR. The counterfactual path explicitly prevents that misattribution.

\subsection{Metric Interpretation and Diagnostic Value}
The all-trial ASR-IFR and conditional cASR-IFR answer complementary questions. ASR-IFR uses the complete paired suite as its denominator and therefore supports stable release-to-release comparison even when some reference-text controls fail. cASR-IFR restricts the denominator to cases in which the frozen downstream system succeeds on reference text and therefore estimates the probability that adding the speech path breaks an otherwise eligible case. In the neural pilot, 142 of 360 trials were speech-attributable, giving 39.4\% all-trial ASR-IFR. Because the text control succeeded in 288 of those trials, the corresponding cASR-IFR was 49.3\%. Reporting both values makes the downstream-control ceiling visible instead of silently excluding failed controls.

This separation also changes the interpretation of lexical error. A high WER trial may remain functionally correct if the downstream system preserves the required intent and entities, while an apparently modest lexical change may redirect the task. The paired examples in Section~V illustrate exactly this distinction: two trials with WER 0.40 had different semantic outcomes. Conversely, an exact transcript is not evidence that the speech system caused a downstream failure; 34 exact-transcript trials occurred in cases whose reference-text controls already failed. CAFA-IVR therefore treats WER as a useful diagnostic dimension rather than as the attribution decision itself.

\subsection{Implications for Regression Testing}
The attribution matrix supports defect routing as well as scoring. A \texttt{SPEECH\_ATTRIBUTABLE} result identifies a regression that appears only after ASR is introduced and is therefore a candidate for speech-layer investigation. A \texttt{DOWNSTREAM\_OR\_TEST} result reproduces without the speech path and should be triaged toward the NLU, agent, workflow, dialogue state, test oracle, or test data. \texttt{CONTEXT\_OR\_ORACLE} exposes the less common case in which audio succeeds while the reference control fails and therefore flags a control inconsistency or context-sensitive behavior for inspection. This routing is operationally distinct from a dashboard that reports only transcript quality and final task success.

For release decisions, CAFA-IVR does not require a universal threshold. Teams can compare a candidate with an approved baseline and combine ASR-IFR with CEER for journeys containing governed entities. Because the reference and speech paths share the same downstream configuration and oracle, the comparison is intended to isolate the incremental effect of adding the speech path. That design does not prove causality outside the controlled test boundary, but it provides a reproducible counterfactual diagnostic for regression engineering.

\section{External Human-Speech Mapping}
The local studies use synthetic speech, so we map the protocol to HarperValleyBank (HVB), a public domain-specific consumer-banking dialogue corpus containing 23.7 hours, 1,446 conversations, 25,730 utterances, 59 speakers, eight task classes, and separate 8-kHz call sides \cite{wu2020hvb}. An online SLUE-HVB sample was frozen across six banking intents with speaker IDs, segment IDs, timestamps, human transcripts, and license/source metadata \cite{shon2022slue2}. This is an external source-validation layer, not a third locally executed ASR experiment: the online audio binaries could not be materialized in the execution environment, so no sample-wide WER or ASR-IFR is claimed. HVB nevertheless supplies the human telephony channel, intent labels, and corrected transcript structure needed by CAFA-IVR without changing the test contract.

\section{Industrial Deployment}
The research contribution is packaged as a versioned reference specification and scorer rather than only a paper. An adopter supplies trial-level outputs from its existing ASR/NLU stack; the scorer produces attributed trials, condition summaries, and baseline/candidate comparisons. A CI job can fail when organization-defined ASR-IFR or CEER regression limits are crossed. Because the same test schema applies to conventional NLU and generative agents, the framework can be inserted around an existing voice stack without replacing its inference components.

\begin{table}[t]
\caption{Minimum adopter interface.}
\label{tab:adopt}
\centering\scriptsize
\begin{tabular}{p{0.26\columnwidth}p{0.64\columnwidth}}
\toprule
Artifact & Purpose\\
\midrule
Reference text & Counterfactual control input\\
Audio + transcript & Speech-path input and lexical diagnostic\\
Task oracle & Frozen success/failure decision for both paths\\
Engine metadata & Reproducible ASR/downstream versioning\\
Baseline summary & Regression comparison in CI/CD\\
\bottomrule
\end{tabular}
\end{table}

Independent implementation should record the CAFA-IVR version, ASR/downstream versions, test corpus, thresholds, resulting engineering decision, and retained evidence. This makes both successful adoption and negative evaluations independently auditable rather than dependent on the original author. The open implementation exposes \texttt{score} and \texttt{compare} commands and adapter examples so a team can adopt the attribution contract without changing its ASR vendor or dialogue platform.

\section{Limitations}
Both locally executed studies use synthetic speech. The neural decoder is closed-set and phrase-overlapping across training/test speakers, so it is an in-domain constrained recognizer rather than a pretrained open-vocabulary benchmark. The 80\% neural text-control accuracy also limits eligible attribution cases. Human-speech multi-engine inference, stronger frozen downstream controls, validated entity normalization, and speaker-level paired statistics are needed before making broader robustness claims.

\section{Conclusion}
CAFA-IVR converts conversational voice regression testing from a single transcript-quality number into a paired attribution procedure. The central contribution is therefore not another transcript-quality score, but an operational test for determining whether the speech/ASR path introduced a downstream failure. Across 540 actual local ASR decodes, it separates lexical mismatch from failures that appear only after speech is introduced and avoids blaming ASR for downstream cases that already fail on reference text. The framework is deliberately implementable as a lightweight layer around existing voice systems, with WER, ASR-IFR, CEER, attribution states, reproducibility metadata, and baseline release gates exposed through a portable specification and reference scorer.

\begin{thebibliography}{00}
\bibitem{kim2021} S. Kim et al., ``Semantic Distance: A New Metric for ASR Performance Analysis Towards Spoken Language Understanding,'' arXiv:2104.02138, 2021.
\bibitem{roy2021} S. Roy, ``Semantic-WER: A Unified Metric for the Evaluation of ASR Transcript for End Usability,'' arXiv:2106.02016, 2021.
\bibitem{sengupta2021} S. Sengupta, J. Krone, and S. Mansour, ``On the Robustness of Intent Classification and Slot Labeling in Goal-oriented Dialog Systems to Real-world Noise,'' in \emph{Proc. NLP4ConvAI}, 2021, pp. 68--79.
\bibitem{awslex} Amazon Web Services, ``Evaluating Amazon Lex V2 bot performance with Test Workbench,'' \emph{Amazon Lex V2 Developer Guide}, accessed Aug. 2026.
\bibitem{gulati2020} A. Gulati et al., ``Conformer: Convolution-augmented Transformer for Speech Recognition,'' in \emph{Proc. Interspeech}, 2020, pp. 5036--5040.
\bibitem{wu2020hvb} M. Wu, J. Nafziger, A. Scodary, and A. Maas, ``HarperValleyBank: A Domain-Specific Spoken Dialog Corpus,'' arXiv:2010.13929, 2020.
\bibitem{shon2022slue2} S. Shon et al., ``SLUE Phase-2: A Benchmark Suite of Diverse Spoken Language Understanding Tasks,'' in \emph{Proc. 61st Annual Meeting of the Association for Computational Linguistics (ACL)}, 2023, pp. 8906--8937.
\end{thebibliography}
\end{document}
